\centerline{\bf \Huge Домашнее задание}
\centerline{\bf \Huge Инверсия}

\begin{flushright}
        \scriptsize{}
\end{flushright}

\problem{1}
Точки $A$ и $B$ лежат на окружности $\omega$. Касательные к окружности, проходящие через точки $A$ и $B$, пересекаются в точке $P$. Докажите, что $P$ является образом середины хорды $A B$ при инверсии относительно $\omega$.

\problem{2}
Рассмотрим вписанную и описанную окружности треугольника $A B C$. Мы уже знаем, что описанная окружность переходит в окружность под действием инверсии относительно вписанной. В какую?

\problem{3}
Окружность $\omega$ переходит под действием инверсии относительно $\Omega$ в себя. Как $\omega$ и $\Omega$ располагаются друг относительно друга?

\problem{4}
Четырёхугольник $A B C D$ вписан в окружность $\Omega$ с центром $O$, причём $O$ не лежит на диагоналях четырёхугольника. Описанная окружность $\Omega_1$ треугольника $A O C$ проходит через середину диагонали $B D$. Докажите, что описанная окружность $\Omega$ треугольника $B O D$ проходит через середину диагонали $A C$.

\problem{5}
Диагонали равнобедренной трапеции пересекаются в точке $P$, а продолжения боковых сторон - в точке $T$. Докажите, что $P$ и $T$ инверсны относительно окружности, диаметром которой служит отрезок, соединяющий середины оснований трапеции, а также относительно окружности, описанной около трапеции.

\problem{6}
Две окружности пересекаются в точках $A$ и $B$. Их общая касательная (та, которая ближе к точке $B$ ) касается окружностей в точках $E$ и $F$. Прямая $A B$ пересекает прямую $E F$ в точке $M$. На продолжении $A M$ за точку $M$ выбрана точка $K$ так, что $K M=M A$. Прямая $K E$ вторично пересекает окружность, содержащую точку $E$, в точке $C$. Прямая $K F$ вторично пересекает окружность, содержащую точку $F$, в точке $D$. Докажите, что точки $C, D$ и $A$ лежат на одной прямой.

\problem{7}
В остроугольном треугольнике $A B C$ проведены высоты $A A_1$ и $B B_1$. Прямая $A_1 B_1$ пересекает описанную окружность треугольника $A B C$ в точках $X$ и $Y$. Окружность, описанная около треугольника $C B B_1$, пересекает высоту $A A_1$ в точке Z. Докажите, что $C X=C Y=C Z$.